%%%%%%%%%%%%%%%%%%%%%%%%%%%%%%%%%%%%%%%%%%%%%%%%%%%%%%%%%%%%%%%%%%%%%%%%%
%
%	LaTeX File for PhD Thesis - Experimental Results
%
%%%%%%%%%%%%%%%%%%%%%%%%%%%%%%%%%%%%%%%%%%%%%%%%%%%%%%%%%%%%%%%%%%%%%%%%%

\chapter{Experimental Results and Analysis}
\label{chap:results}

\section{Introduction}
\label{res:intro}
This chapter presents a comprehensive evaluation of the methods developed in this thesis. We consolidate the experimental results for DenoMAE (Chapter 4), Plug-and-Play LLM (Chapter 5), and Meta-Task (Chapter 6) to provide a holistic analysis of their performance. The goal is to validate the effectiveness of each contribution and demonstrate how they collectively form a robust framework for wireless signal analysis.

\section{Evaluation Methodology}
\label{res:eval}
This section summarizes the experimental setup, datasets, and evaluation metrics used to validate our proposed methods. A more detailed description of the overall methodology is provided in Chapter 3.

\paragraph{Datasets}
We use the RadioML 2016.10a dataset for both pre-training and fine-tuning. For pre-training, we use the full unlabeled dataset. For fine-tuning, we use a small fraction of the labeled data, with varying numbers of samples per class to evaluate the data efficiency of our approach.

\paragraph{Evaluation Metrics}
We evaluate the performance of our models using the following metrics:
\begin{itemize}
    \item \textbf{Classification Accuracy:} The primary metric for evaluating the performance of our models.
    \item \textbf{SNR-Specific Performance:} We analyze the performance of our models across a range of signal-to-noise ratios (SNRs) to assess their robustness to noise.
    \item \textbf{Confusion Matrices:} We use confusion matrices to analyze the classification performance for each modulation class.
    \item \textbf{Statistical Significance Testing:} We use statistical tests to determine if the performance differences between models are statistically significant.
\end{itemize}

\section{DenoMAE Results}
\label{res:denomae}
\subsection{Overall Classification Performance}
Table~\ref{tab:denomae_snr_performance} presents DenoMAE's classification accuracy across SNR conditions compared to baseline methods. DenoMAE achieves 83.5\% overall accuracy with only 10K unlabeled pretraining samples and 1K labeled fine-tuning samples, representing 10× and 3× data reduction respectively.

\begin{table}[ht]
\centering
\caption{DenoMAE classification accuracy (\%) across SNR levels. DenoMAE (Ours) vs. non-pretrained baseline and state-of-the-art methods.}
\label{tab:denomae_snr_performance}
\begin{tabular}{lcccccc}
\toprule
\textbf{Method} & \textbf{-10 dB} & \textbf{-5 dB} & \textbf{0 dB} & \textbf{+5 dB} & \textbf{+10 dB} & \textbf{Avg.} \\
\midrule
Non-pretrained ViT & 45.2 & 61.3 & 78.1 & 87.4 & 92.1 & 72.8 \\
CNN Baseline & 48.7 & 64.5 & 79.8 & 88.2 & 92.8 & 74.8 \\
Transformer (Kong) & 51.3 & 66.9 & 81.2 & 89.1 & 93.4 & 76.4 \\
\midrule
\textbf{DenoMAE (Ours)} & \textbf{67.3} & \textbf{79.8} & \textbf{88.7} & \textbf{93.6} & \textbf{95.8} & \textbf{85.0} \\
\textbf{Improvement} & \textbf{+22.1} & \textbf{+18.5} & \textbf{+10.6} & \textbf{+6.2} & \textbf{+3.7} & \textbf{+12.2} \\
\bottomrule
\end{tabular}
\end{table}

\textbf{Key Findings:} DenoMAE's advantage is most pronounced in low-SNR conditions (+22.1\% at -10 dB), validating that treating noise as an informative modality enhances robustness. Performance gains diminish at high SNR where noise is less influential.

\subsection{SNR Extrapolation}
DenoMAE generalizes to unseen SNR conditions. When pretrained on [-10, +10] dB and tested on [-20, -10] dB, accuracy remains above 55\%, demonstrating learned signal structure transfers beyond training distribution.

\subsection{Ablation Studies}
Table~\ref{tab:denomae_ablation} evaluates the contribution of each DenoMAE component. Both multimodal processing and noise-as-modality are essential for optimal performance.

\begin{table}[ht]
\centering
\caption{DenoMAE ablation study. Contribution of each component to overall accuracy.}
\label{tab:denomae_ablation}
\begin{tabular}{lcc}
\toprule
\textbf{Configuration} & \textbf{Low SNR (-10 to 0 dB)} & \textbf{Overall Acc.} \\
\midrule
Single modality (IQ only) & 58.3 & 78.4 \\
Single modality (Constellation only) & 61.7 & 80.1 \\
Multimodal (no noise modeling) & 64.2 & 82.3 \\
\textbf{Full DenoMAE (Ours)} & \textbf{78.6} & \textbf{85.0} \\
\bottomrule
\end{tabular}
\end{table}

\subsection{Data Efficiency Analysis}
Table~\ref{tab:denomae_data_efficiency} demonstrates DenoMAE's data efficiency. With only 1K labeled samples (vs. 10K for supervised baseline), DenoMAE achieves 83.5\% accuracy, matching supervised performance with 10× less labeled data.

\begin{table}[ht]
\centering
\caption{Data efficiency comparison. DenoMAE vs. fully supervised training.}
\label{tab:denomae_data_efficiency}
\begin{tabular}{lccc}
\toprule
\textbf{Method} & \textbf{Unlabeled Pretrain} & \textbf{Labeled Finetune} & \textbf{Accuracy} \\
\midrule
Supervised Baseline & 0 & 10,000 & 83.2 \\
Supervised Baseline & 0 & 3,000 & 78.6 \\
Supervised Baseline & 0 & 1,000 & 72.8 \\
\midrule
\textbf{DenoMAE} & 10,000 & 1,000 & \textbf{83.5} \\
\textbf{DenoMAE} & 10,000 & 500 & \textbf{81.2} \\
\textbf{DenoMAE} & 10,000 & 100 & \textbf{76.4} \\
\bottomrule
\end{tabular}
\end{table}

\section{Plug-and-Play LLM Results}
\label{res:pnp}

\subsection{Classification Accuracy}
Table~\ref{tab:pnp_performance} compares Plug-and-Play performance across different LLM architectures. Training-free classification achieves 72.4\% accuracy with o3-mini, competitive with extensively trained models while requiring zero wireless-specific training.

\begin{table}[ht]
\centering
\caption{Plug-and-Play LLM classification accuracy. Training-free vs. trained baselines.}
\label{tab:pnp_performance}
\begin{tabular}{lccc}
\toprule
\textbf{Method} & \textbf{Parameters} & \textbf{Training Cost} & \textbf{Accuracy (\%)} \\
\midrule
\multicolumn{4}{l}{\textit{Trained Baselines:}} \\
CNN Supervised & 2.3M & Full training & 74.8 \\
Transformer Supervised & 12M & Full training & 76.4 \\
DenoMAE (pretrain+finetune) & 86M & Self-supervised & 85.0 \\
\midrule
\multicolumn{4}{l}{\textit{Training-Free LLM (Ours):}} \\
DeepSeek-R1-Distill-Qwen-7B & 7B & \textbf{Zero} & 68.2 \\
DeepSeek-R1-Distill-Qwen-32B & 32B & \textbf{Zero} & 70.5 \\
\textbf{OpenAI o3-mini (200B)} & 200B & \textbf{Zero} & \textbf{72.4} \\
\bottomrule
\end{tabular}
\end{table}

\textbf{Key Insight:} While 72.4\% is lower than DenoMAE's 85\%, the zero-training deployment enables immediate use in scenarios where labeled data is unavailable or time-to-deployment is critical.

\subsection{SNR Robustness Analysis}
Table~\ref{tab:pnp_snr} shows Plug-and-Play performance across noise conditions. Training-free classification maintains 65-70\% accuracy even at low SNR, demonstrating robustness without noise-specific training.

\begin{table}[ht]
\centering
\caption{Plug-and-Play accuracy across SNR levels (o3-mini 200B).}
\label{tab:pnp_snr}
\begin{tabular}{lcccc}
\toprule
\textbf{SNR Range} & \textbf{Accuracy (\%)} & \textbf{vs. DenoMAE} & \textbf{Training Required} \\
\midrule
-10 to 0 dB (Low) & 65.3 & -13.3 & None \\
0 to 10 dB (Medium) & 72.8 & -15.9 & None \\
10 to 20 dB (High) & 76.1 & -19.5 & None \\
\midrule
\textbf{Overall} & \textbf{72.4} & \textbf{-12.6} & \textbf{Zero} \\
\bottomrule
\end{tabular}
\end{table}

\textbf{Trade-off Analysis:} Plug-and-Play sacrifices 12.6\% accuracy for complete training elimination. For rapid prototyping, emergency deployment, or exploratory analysis, this trade-off is highly favorable.

\subsection{Prompt Engineering Impact}
Statistical feature selection critically impacts performance. Table~\ref{tab:pnp_features} shows ablation on prompt components.

\begin{table}[ht]
\centering
\caption{Ablation on statistical features in prompts.}
\label{tab:pnp_features}
\begin{tabular}{lc}
\toprule
\textbf{Feature Set} & \textbf{Accuracy (\%)} \\
\midrule
Basic stats only (mean, variance) & 58.2 \\
+ Higher-order moments (up to 7th) & 64.7 \\
+ k-statistics (cumulants 1-4) & 69.8 \\
\textbf{Full feature set (Ours)} & \textbf{72.4} \\
\bottomrule
\end{tabular}
\end{table}

\section{Meta-Task Results}
\label{res:metatask}
\subsection{Few-Shot Learning Performance}
Our Meta-Task framework consistently improves the performance of all baseline models across all datasets and settings. For example, on the Mini-ImageNet dataset, our framework improves the accuracy of Prototypical Networks by 2.5\% in the 5-way 1-shot setting and by 1.8\% in the 5-way 5-shot setting. Similarly, for MAML, we observe an improvement of 2.1\% and 1.5\% in the 1-shot and 5-shot settings, respectively.

The improvements are even more significant on the more challenging Tiered-ImageNet and FC100 datasets. On Tiered-ImageNet, our framework improves the accuracy of MetaOptNet by 3.2\% in the 1-shot setting and by 2.5\% in the 5-shot setting. On FC100, we observe an improvement of 3.5\% and 2.8\% for Prototypical Networks in the 1-shot and 5-shot settings, respectively.

\subsection{Adaptation Speed Analysis}
The Meta-Task framework demonstrates faster convergence compared to baseline models, as shown by the learning curves in the original paper.

\subsection{Cross-Domain Generalization}
The framework's ability to generalize to new tasks is a key advantage. The consistent performance improvements across different datasets highlight its cross-domain generalization capabilities.

\subsection{Ablation Study}
To understand the contribution of each component of our framework, we conduct an ablation study on the Mini-ImageNet dataset. We evaluate the performance of our framework with and without the Task-Decoder. The results are shown in Table \ref{tab:ablation_study_results}.

As expected, removing the Task-Decoder leads to a significant drop in performance. For example, for Prototypical Networks, the accuracy drops by 2.5\% in the 1-shot setting and by 1.8\% in the 5-shot setting. This demonstrates the importance of the Task-Decoder in improving the generalization performance of the model.

\begin{table}[ht]
\centering
\caption{Ablation study on Mini-ImageNet. We report the mean classification accuracy (\%) with 95\% confidence intervals.}
\label{tab:ablation_study_results}
\begin{tabular}{lcc}
\toprule
\textbf{Method} & \textbf{5-way 1-shot} & \textbf{5-way 5-shot} \\
\midrule
Prototypical Networks & 49.42 $\pm$ 0.78 & 68.20 $\pm$ 0.66 \\
- Task-Decoder & 46.92 $\pm$ 0.79 & 66.40 $\pm$ 0.68 \\
\midrule
MAML & 48.70 $\pm$ 0.84 & 63.11 $\pm$ 0.92 \\
- Task-Decoder & 46.60 $\pm$ 0.85 & 61.61 $\pm$ 0.94 \\
\bottomrule
\end{tabular}
\end{table}

\section{Synthesizing the Modular Framework}
\label{res:synth}
The experimental results from the preceding sections~\ref{res:denomae}, \ref{res:pnp}, \ref{res:metatask} validate the three core components of our proposed framework, each addressing a distinct challenge in developing cost-effective wireless foundation models. This section synthesizes these individual findings to demonstrate how they form the cohesive, modular system proposed by this dissertation.

\textbf{1. The Representation Layer (Validated by DenoMAE, Section~\ref{res:denomae}):}
The framework's foundation is an efficient and robust feature encoder. Section~\ref{res:denomae} validated this component through DenoMAE. The key finding—achieving 85.0\% accuracy while using 10x less labeled data than a supervised baseline—proves that robust, universal signal representations can be learned efficiently from unlabeled data. This component successfully provides the high-quality feature backbone that any modular system requires, validating the representation layer.

\textbf{2. The Reasoning Layer (Validated by Plug-and-Play LLM, Section~\ref{res:pnp}):}
A key capability of a foundation model is reasoning and generalization. Section~\ref{res:pnp} validated this by demonstrating a training-free deployment path. The 72.4\% accuracy achieved by the Plug-and-Play LLM, which required zero wireless-specific training, proves that pre-trained world knowledge from LLMs can be directly transferred to the wireless domain. This provides an immediate, cost-free reasoning layer for classification, interpretation, and open-set analysis.

\textbf{3. The Adaptation Layer (Validated by Meta-Task, Section~\ref{res:metatask}):}
Finally, a foundation model must be adaptable to new, unforeseen tasks. Section~\ref{res:metatask} validated this adaptation layer. The consistent accuracy improvements (e.g., +2.5\% in 1-shot, +5.9\% in 5-shot) and faster convergence (20k vs. 50k episodes) demonstrate that the framework can be rapidly specialized to new tasks with minimal data. This addresses the critical need for customization without costly, full-scale retraining.

In summary, while a full end-to-end experimental integration of these components is a key direction for future work (as discussed in Section 7.8), the individual validation of each module confirms their standalone viability and complementary roles. The results prove that the proposed framework is not just a theoretical construct but a modular system where the representation, reasoning, and adaptation components are individually proven effective and ready for synergistic deployment.


\subsection{Answering the Research Questions}

\subsubsection{RQ1: Can Vision-Based Self-Supervised Learning Transfer World Knowledge to Wireless?}

\textbf{Answer: Yes, with proper signal-to-image representations.}

Our work provides strong evidence through DenoMAE (Chapter~\ref{chap:denomae}) and CDRF validation (Chapter~\ref{chap:cdrf}):

\textbf{DenoMAE Findings:}
\begin{itemize}
    \item Constellation diagrams serve as effective visual representations, enabling Vision Transformers to process wireless signals
    \item Masked autoencoding learns universal signal features with 80\% reduction in labeled data requirements
    \item Multimodal learning (IQ samples + constellation diagrams) improves robustness across SNR conditions
    \item Features learned transfer effectively to modulation classification and signal detection tasks
\end{itemize}

\textbf{CDRF Cross-Domain Validation:}
\begin{itemize}
    \item YOLO architecture successfully detects RF drones using spectrogram representations
    \item Vision models' pre-trained knowledge transfers to wireless domain without wireless-specific training
    \item Demonstrates bidirectional understanding: signal models understand images, vision models understand signals
\end{itemize}

\textbf{Key Insight:} The signal-to-image transformation (constellation diagrams, spectrograms) acts as a bridge, allowing pre-trained vision models to approximate world knowledge for wireless applications. This cross-modal approach dramatically reduces dependency on large labeled wireless datasets.

\subsubsection{RQ2: Can LLMs Enable Open-Set Classification Through World Knowledge Transfer?}

\textbf{Answer: Yes, competitive with extensively trained models.}

The Plug-and-Play framework (Chapter~\ref{chap:plugandplay}) demonstrates:

\textbf{Technical Achievements:}
\begin{itemize}
    \item Statistical signal features (higher-order cumulants, spectral properties) successfully convert to natural language prompts
    \item Pre-trained LLMs (DeepSeek-R1, OpenAI o3-mini) achieve 72.4\% accuracy without wireless-specific fine-tuning
    \item Training-free operation enables rapid deployment without expensive data collection or model training
    \item LLMs provide interpretable, human-understandable explanations for classification decisions
\end{itemize}

\textbf{Semantic Gap Bridging:} The key innovation lies in the prompt template design that translates quantitative signal properties into linguistic descriptions LLMs can reason about. This demonstrates that sophisticated reasoning capabilities need not be trained from scratch but can leverage existing world knowledge.

\textbf{Competitive Performance:} While 72.4\% is lower than specialized deep learning models (83.5\% for DenoMAE), the training-free nature represents a fundamentally different trade-off. For scenarios where labeled data is unavailable or deployment speed is critical, LLM-based classification provides immediate capability.

\subsubsection{RQ3: Can Meta-Learning Enable Rapid Adaptation with Minimal Data?}

\textbf{Answer: Yes, with significant performance improvements.}

Meta-Task (Chapter~\ref{chap:metatask}) validates rapid adaptation:

\textbf{Quantitative Results:}
\begin{itemize}
    \item 4.3 percentage point accuracy improvement over baselines with only 5-10 labeled examples per class
    \item Faster convergence (20k episodes vs 50k for standard training)
    \item Method-agnostic design validated across Prototypical Networks, MAML, MetaOptNet, and P>M>F
    \item Cross-dataset generalization demonstrates learned adaptation strategies transfer
\end{itemize}

\textbf{Design Principles Validated:}
\begin{itemize}
    \item Episodic training with SNR-aware episode generation captures wireless-specific characteristics
    \item Task-aware regularization (Task-Decoder) prevents overfitting while maintaining adaptability
    \item Learnable regularization as auxiliary task improves generalization without hyperparameter tuning
\end{itemize}

\textbf{Practical Impact:} The ability to adapt with 5-10 examples addresses the cost barrier for wireless foundation model deployment. Rather than requiring thousands of labeled samples for each new channel condition or modulation scheme, Meta-Task enables economical adaptation.

\subsection{Interpretation of Results}

\subsubsection{Why DenoMAE Succeeds}

DenoMAE's effectiveness stems from three synergistic design choices:

\begin{enumerate}
    \item \textbf{Multimodal Learning:} Joint processing of time-domain IQ samples and constellation diagrams creates richer representations than single modalities. The model learns complementary features—temporal dynamics from IQ samples, geometric structure from constellation diagrams.
    
    \item \textbf{Noise as Informative Modality:} Rather than treating noise as corruption to remove, DenoMAE models it explicitly. This forces the encoder to learn signal structure independent of noise characteristics, leading to superior low-SNR performance.
    
    \item \textbf{Vision Transformer Architecture:} The global receptive field and self-attention mechanism naturally capture long-range dependencies in signals. Pre-training strategies (masked autoencoding) proven in computer vision transfer effectively to wireless domain when signals are represented visually.
\end{enumerate}

\subsubsection{LLM Reasoning Advantages}

The Plug-and-Play framework represents a paradigm shift from black-box classifiers to interpretable reasoning systems:

\begin{enumerate}
    \item \textbf{Training-Free Deployment:} Eliminates data collection, labeling, and training costs. For rapid prototyping or scenarios where labeled data is unavailable, this provides immediate classification capability.
    
    \item \textbf{Interpretability:} Natural language reasoning provides transparent decision-making. Engineers can understand why the model classified a signal a certain way, enabling debugging and refinement.
    
    \item \textbf{Open-Set Recognition Potential:} LLMs' vast world knowledge provides a foundation for handling unknown signal types. While not fully realized in current work (closed-set evaluation), the architecture naturally extends to open-set scenarios.
\end{enumerate}

\subsubsection{Meta-Learning Benefits and Synergies}

Meta-Task's success demonstrates that adaptation can itself be learned:

\begin{enumerate}
    \item \textbf{Rapid Adaptation:} The framework learns how to adapt, not just what to adapt to. This meta-level learning enables quick specialization for new tasks with minimal data.
    
    \item \textbf{Method-Agnostic Design:} Compatibility with multiple meta-learning algorithms (Prototypical Networks, MAML, MetaOptNet) demonstrates that learnable regularization is a general principle, not algorithm-specific.
    
    \item \textbf{Wireless-Aware Episodes:} SNR-stratified episode generation accounts for signal-specific characteristics, improving generalization across channel conditions.
\end{enumerate}

\textbf{Framework Synergies:} The three components work synergistically: (1) DenoMAE provides rich, transferable representations; (2) Meta-Task enables rapid adaptation of these representations; (3) Plug-and-Play offers training-free reasoning when adaptation data is unavailable. Together, they address complementary aspects of wireless foundation model development.

% \section{Computational Efficiency Analysis}
% We analyze the computational cost of our framework, including training time, inference time, and memory requirements, demonstrating its efficiency compared to traditional methods.

% \section{Summary of Findings}
% The experimental results presented in this chapter validate the effectiveness of the proposed methods. DenoMAE establishes a strong foundation with its data-efficient and robust representation learning. The Plug-and-Play LLM framework demonstrates the potential for training-free, interpretable classification. Finally, the Meta-Task framework showcases the ability to rapidly adapt to new tasks with minimal data. Together, these results support the viability of a multi-stage framework for building powerful and flexible wireless foundation models.


% \subsection{Contribution 1: DenoMAE - Multimodal Self-Supervised Representation Learning}
% \label{subsec:denomae_contribution}

% \textbf{Addresses RQ1 and RQ5:} Can multimodal self-supervised learning with image representations approximate world knowledge? Is there bidirectional understanding between signal and image domains?

% DenoMAE introduces a novel self-supervised learning framework that learns universal wireless signal representations from unlabeled data through multimodal masked autoencoders:

% \begin{itemize}
%     \item \textbf{Multimodal Processing:} Jointly processes time-domain IQ samples and constellation diagrams (as images) through modality-specific encoders, leveraging vision transformer architectures
%     \item \textbf{Noise as Informative Modality:} Treats noise explicitly as an informative modality rather than corruption, learning to reconstruct both signal content and noise characteristics
%     \item \textbf{Masked Reconstruction:} Masks 75\% of patches across modalities and reconstructs complete signals, forcing the model to learn rich representations of signal structure
%     \item \textbf{Dramatic Data Efficiency:} Achieves 83.5\% accuracy with only 10,000 unlabeled pretraining samples and 1,000 labeled fine-tuning samples (10× reduction in pretraining data, 3× reduction in labeled data)
%     \item \textbf{SNR Robustness:} Demonstrates robust performance from -10 dB to +10 dB, with up to 22.1\% accuracy improvement over non-pretrained models at -10 dB
%     \item \textbf{Extrapolation Capability:} Successfully generalizes to unseen lower SNRs (down to -20 dB), demonstrating true understanding of signal properties
% \end{itemize}

% \textit{Key Innovation:} DenoMAE validates that vision transformer architectures can approximate world knowledge when wireless signals are converted to image representations (constellation diagrams), demonstrating bidirectional understanding between signal and image domains. This enables transfer learning from computer vision while maintaining wireless-specific signal understanding.

% \textit{Significance:} Provides the rich feature encoder component for wireless foundation models, demonstrating that universal representations can be learned economically through self-supervised learning without massive labeled datasets. Addresses the fundamental data scarcity challenge in wireless AI.

% \subsection{Contribution 2: Plug-and-Play - LLM-Based Training-Free Classification}
% \label{subsec:plugandplay_contribution}

% \textbf{Addresses RQ2:} Can LLMs enable open-set classification of wireless signals through world knowledge transfer?

% The Plug-and-Play framework introduces a pioneering proof-of-concept that enables training-free wireless classification through Large Language Model reasoning:

% \begin{itemize}
%     \item \textbf{Training-Free Operation:} Utilizes pre-trained LLMs (DeepSeek-R1-Distill-Qwen 7B/32B, OpenAI o3-mini 200B) to perform wireless classification without any wireless-specific training
%     \item \textbf{Statistical Feature-to-Language Conversion:} Converts quantitative signal features (higher-order cumulants, spectral characteristics) into structured natural language prompts that LLMs can reason about
%     \item \textbf{In-Context Learning:} Leverages LLMs' inherent in-context learning (ICL) capabilities for one-shot classification, providing exemplar context (I+C+S prompts) that substantially improves performance
%     \item \textbf{Competitive Performance:} Achieves up to 72.4\% cleaned accuracy under noisy conditions with o3-mini, competitive with extensively trained task-specific models
%     \item \textbf{Open-Set Capability:} Enables reasoning about signal types beyond traditional closed-set classification through natural language understanding
%     \item \textbf{Interpretability:} Provides human-understandable explanations for classification decisions, facilitating prompt-based optimization and interaction
% \end{itemize}

% \textit{Key Innovation:} Plug-and-Play bridges the semantic gap between wireless signal statistics and linguistic descriptions, demonstrating that sophisticated reasoning need not be trained from scratch but can leverage pre-trained world knowledge from LLMs. This validates cross-domain knowledge transfer from language to wireless domains.

% \textit{Significance:} This proof-of-concept validates the feasibility of integrating LLMs as intelligent reasoning engines for wireless applications, drastically reducing deployment costs while providing interpretability—a critical requirement for wireless systems. Represents a novel approach to achieving foundation model reasoning capabilities economically.

% \subsection{Contribution 3: Meta-Task - Rapid Adaptation through Meta-Learning}
% \label{subsec:metatask_contribution}

% \textbf{Addresses RQ3:} Can meta-learning enable rapid adaptation to new tasks with minimal labeled data?

% Meta-Task introduces task-aware regularization for meta-learning, enabling rapid adaptation to new wireless classification tasks:

% \begin{itemize}
%     \item \textbf{Task-Aware Regularization:} Reformulates regularization as a learnable task within the few-shot learning optimization process, introducing parameterized neural layers that dynamically learn task-specific regularization strategies
%     \item \textbf{Task-Decoder Architecture:} Employs unsupervised autoencoders to reconstruct input representations from embeddings, refining hidden representations by constraining accurate reconstruction
%     \item \textbf{Method-Agnostic Design:} Seamlessly integrates with multiple meta-learning algorithms (Prototypical Networks, MAML, MetaOptNet, P>M>F), demonstrating broad applicability
%     \item \textbf{Significant Performance Gains:} Improves Prototypical Networks from 46.6\% to 50.0\% (5-way 1-shot) and 65.4\% to 71.3\% (5-way 5-shot) on benchmarks like MiniImageNet and TieredImageNet
%     \item \textbf{Faster Convergence:} Achieves higher accuracy within 20k episodes compared to baseline methods requiring 50k episodes
%     \item \textbf{Minimal Overhead:} Maintains computational efficiency through shallow decoder architectures (75K-250K parameters)
% \end{itemize}

% \textit{Key Innovation:} Meta-Task addresses the fundamental overfitting challenge in meta-learning by learning regularization strategies specific to each task, rather than applying fixed regularization. This enables better generalization across diverse wireless tasks and channel conditions.

% \textit{Significance:} Validates that meta-learning can enable rapid adaptation to new wireless classification tasks with minimal labeled data (5-10 examples), providing an efficient adaptation mechanism essential for cost-effective foundation model deployment to new scenarios.

% \subsection{Contribution 4: CDRF - A Benchmark for Cross-Modal RF Perception}
% \label{subsec:cdrf_contribution}

% \textbf{Addresses RQ4:} Can vision models transfer world knowledge to wireless domains through cross-modal learning?

% We introduce the CageDroneRF (CDRF) benchmark, a large-scale dataset and toolkit for RF-based drone detection, classification, and open-set recognition. This contribution provides a critical resource and methodology for the community:

% \begin{itemize}
%     \item \textbf{Large-Scale RF Dataset:} Comprises over 500 GB of real-world RF captures from 39 classes, including diverse drone models, controllers, and ambient noise, recorded in both controlled RF-cage and outdoor environments.
%     \item \textbf{Raw-Signal Augmentation Pipeline:} Provides open-source tools for generating realistic training data by controlling SNR, injecting interference, and applying frequency shifts at the raw I/Q signal level.
%     \item \textbf{Cross-Modal Validation:} Demonstrates that by converting RF signals into spectrogram images, state-of-the-art computer vision object detectors (e.g., YOLO) can be successfully applied to RF perception tasks like drone detection.
%     \item \textbf{Standardized Evaluation:} Enables rigorous, reproducible benchmarking for detection, classification, and open-set recognition, accelerating progress toward robust, generalizable RF models.
% \end{itemize}

% \textit{Key Innovation:} CDRF provides the first comprehensive, open-source benchmark with an integrated toolkit for raw-signal augmentation and evaluation. It validates a practical methodology for transferring knowledge from pre-trained vision models to the wireless domain.

% \textit{Significance:} This work provides the research community with a vital resource for developing and testing robust drone detection systems. It empirically answers RQ4 by showing that vision models can indeed transfer learned spatial feature hierarchies to wireless tasks, paving the way for new cross-modal applications.


% \section{Summary of Contributions and Answers to Research Questions}

% This section summarizes how the four core contributions of this dissertation collectively answer the five research questions, establishing a modular and cost-effective pathway toward Wireless Physical-Layer Foundation Models (WPFMs).

% \begin{itemize}
%     \item \textbf{DenoMAE (Chapter~\ref{chap:denomae}):} Addresses RQ1 and RQ5 by demonstrating that multimodal self-supervised learning with image representations can approximate world knowledge, achieving a 10× reduction in pretraining data and a 3× reduction in labeled data while enabling bidirectional understanding between signal and image domains.
    
%     \item \textbf{CDRF (Chapter~\ref{chap:cdrf}):} Addresses RQ4 by providing a large-scale benchmark and toolkit that validates the transfer of knowledge from computer vision models to RF perception tasks, demonstrating a practical methodology for cross-modal learning.

%     \item \textbf{Plug-and-Play (Chapter~\ref{chap:plugandplay}):} Addresses RQ2 by validating that LLMs can enable training-free, open-set wireless classification through statistical feature-to-language conversion, achieving competitive accuracy (72.4\%) without wireless-specific training.

%     \item \textbf{Meta-Task (Chapter~\ref{chap:meta_task}):} Addresses RQ3 by enabling rapid adaptation to new tasks with 5-10 labeled examples through task-aware regularization, improving few-shot accuracy by up to 4.3 percentage points.
% \end{itemize}
